% Methods — see paper/methods_skeleton.md for the long-form outline.
% Target 1500-2000 words for PLOS Comp Bio. Each subsection below maps to
% a markdown subsection in methods_skeleton.md.

\subsection{Model overview}\label{sec:methods:overview}
The two-leg spinal CPG presented here extends a canonical computational
sCPG model~\cite{rybak2006,shevtsova2015} with four
mechanisms not present in the Zhang 2022 reference model~\cite{zhang2022}:
(i) STDP-driven self-organisation of descending and cutaneous afferent
synapses onto the rhythm generator;
(ii) a closed-loop Ia feedback pathway to the reciprocal-inhibition
interneurons;
(iii) motoneuron pools driving Hill-like muscle proxies with force and
length dynamics;
(iv) externally paced sequential Ia activation simulating heel-to-toe
weight transfer during stance.
The model is implemented in NEST~3.9~\cite{nest2002} and is publicly
available at \href{https://github.com/max-talanov/tinyCPG}{github.com/max-talanov/tinyCPG}.

\begin{figure}[t]
  \centering
  \includegraphics[width=0.8\textwidth]{fig1_schematic.png}
  \caption{\textbf{Architecture of the two-leg spinal CPG.}
    Each leg contains a flexor (F)
    and extensor (E) rhythm-generator half-centre coupled by asymmetric
    reciprocal inhibition via InE and InF interneurons (Zhang 2022 6:1
    ratio). Sensory afferents (Ia, cutaneous) feed into the
    reciprocal-inhibition interneurons; motor pools (M-E, M-F) drive
    muscle proxies whose force and length signals close the Ia loop.
    L\,$\leftrightarrow$\,R commissural inhibition (V0-like) is mediated by
    flexor- and extensor-targeted projections.}
  \label{fig:schematic}
\end{figure}

\subsection{Single-neuron model}\label{sec:methods:neuron}
The network is implemented at two complementary levels of biophysical
detail, using the same architecture, the same plastic projections and the
same experimental protocol in both cases.

\paragraph{Reduced (Izhikevich) implementation.}
In the NEST implementation, all spiking neurons are modelled as Izhikevich
quadratic-integrate-fire units~\cite{izhikevich2003}. Each neuron obeys
\begin{align}
  \frac{dv}{dt} &= 0.04v^2 + 5v + 140 - u + I, \\
  \frac{du}{dt} &= a\,(bv - u),
\end{align}
with the after-spike reset $v \leftarrow c$, $u \leftarrow u + d$ applied
whenever $v \geq 30$\,mV. Excitatory extensor RG-E neurons use the
regular-spiking set ($a=0.02, b=0.2, c=-65, d=8$). Flexor RG-F neurons use
the intrinsically-bursting set ($a=0.02, b=0.2, c=-55, d=4$), which
approximates persistent-sodium-driven ($I_\text{NaP}$) bursting
phenomenologically rather than deriving it from channel kinetics.
Inhibitory interneurons (InE, InF, Ia-INT) use a fast-spiking set
($d=2$). This formulation reproduces the spiking and bursting regimes
required by the rhythm generator at a small fraction of the computational
cost of a conductance-based model, which is what makes the large
mode\,$\times$\,learning-rate parameter sweeps reported here tractable.

\paragraph{Conductance-based (Hodgkin--Huxley) implementation.}
Because the reduced model represents bursting phenomenologically, we also
implemented the same circuit with conductance-based neurons in
NEURON~\cite{hines1997}, following the Hodgkin--Huxley
formalism~\cite{hodgkinHuxley1952}. Membrane potential evolves as
\begin{equation}
  C_m \frac{dV}{dt} = -\!\sum_i g_i\, m_i^{p}h_i^{q}\,(V - E_i) - I_\text{leak} + I_\text{syn},
  \label{eq:hh}
\end{equation}
where each active conductance $g_i$ is gated by activation ($m_i$) and
inactivation ($h_i$) variables obeying first-order kinetics
$\dot{x} = (x_\infty(V) - x)/\tau_x(V)$. Rhythm-generator neurons include
fast sodium, delayed-rectifier potassium and leak conductances together
with an explicit persistent sodium current ($I_\text{NaP}$), so that
bursting arises from channel kinetics rather than from a reset rule.
Motoneurons and interneurons are multi-compartment cells with somatic and
dendritic sections. This implementation resolves the sub-threshold and
channel-level dynamics that the Izhikevich formulation abstracts away, at
substantially higher computational cost; it therefore serves as a
biophysical cross-check on the reduced model rather than as the vehicle
for the full parameter sweep.%
\todo{The NEURON/HH results are not yet included in Results. Until they
are, this subsection describes an implementation the paper does not yet
report on, and the Limitations no longer flag the absence of a
conductance-based model --- add the HH results (or re-scope this
paragraph) before circulating.}

Synaptic delays in both implementations are drawn from a rat
\texttt{length\_velocity} preset (Table~\ref{tab:delays}) with
$\sigma=0.2$\,ms Gaussian jitter.

\subsection{Network architecture}\label{sec:methods:network}

\paragraph{Asymmetric reciprocal inhibition (Zhang 2022).}
The F$\to$InF$\to$E pathway is strong ($W_{\text{InF}\to\text{RG-E}}=-48$\,pA,
$P=0.30$); the E$\to$InE$\to$F pathway is weak
($W_{\text{InE}\to\text{RG-F}}=-8$\,pA, $P=0.15$).
The 6:1 ratio preserves the flexor's role as the rhythm-leading
half-centre~\cite{zhang2022,talpalar2013}.

\paragraph{Closed-loop Ia$\to$reciprocal-IN pathway (this work).}
Ia-E parrots project to InE ($W=6$\,pA, $P=0.25$) and Ia-F parrots to
InF, providing closed-loop sensory drive into the reciprocal-inhibition
core. This pathway compensates the loss of strong tonic descending drive
characteristic of deafferented fictive locomotion~\cite{pearson1995,
rossignol2006,hultborn2006}.

\paragraph{Cutaneous reflex.}
Cutaneous parrots project to RG-E (STDP-plastic) and to InE
($W=6$\,pA, $P=0.30$), implementing the stance-phase cutaneous-reflex
gating of extensor activity~\cite{forssberg1980,pearson1995}.

\paragraph{Commissural inhibition.}
L$\leftrightarrow$R flexor RG inhibition is strong
($W=-20$\,pA, $P=0.22$); extensor inhibition is weak
($W=-8$\,pA, $P=0.10$). This V0-mediated commissural inhibition
enforces left--right alternation~\cite{kiehn2016,talpalar2013}.

\paragraph{Motor pool, muscle relay.}
Each motor neuron projects to $N_{\text{mus}}=100$ muscle-parrot relays
($P_{\text{M}\to\text{mus}}=0.8$). Motor pool reciprocal inhibition
($W=-22$\,pA, $P=0.25$) enforces antagonist suppression at the motor
level~\cite{mccreaRybak2008}.

\paragraph{Connectivity specification.}
The full per-connection weight and delay distributions of every
projection were obtained by querying NEST directly after network
construction. Table~\ref{tab:delays}
lists all projections grouped by the functional blocks of the
architecture (Fig.~\ref{fig:schematic}): descending/supraspinal drive,
the sensory afferent$\to$RG projections, the rhythm-generator reciprocal
core, the Ia proprioceptive interneuron pathway, motor output, and the
commissural (interlimb) projections. Static projections are
heterogeneous: their per-connection weights are drawn from a lognormal
distribution about the design value (coefficient of variation $0.5$, mean
and sign preserved), since biological synaptic weights are lognormally
distributed rather than uniform~\cite{song2005,buzsaki2014}. The mean
preserves the designed balance---in particular the asymmetric reciprocal
inhibition (InF$\to$RG-E $=-48$\,pA, InE$\to$RG-F $=-8$\,pA) reproduces the
Zhang~6:1 ratio exactly. The plastic projections (CUT$\to$RG-E in both
arms, BS$\to$RG in the descending arm, Ia$\to$RG in the sensory arm) carry
lognormal weight distributions from STDP: the table reports their
\emph{converged} mean$\pm$s.d.\ (CUT$\to$RG-E and Ia$\to$RG at their
learned load set-points, BS$\to$RG frozen at weak init in the sensory
arm). The cutaneous input to RG-E is a single plastic projection. Conduction delays follow the rat
\texttt{length\_velocity} preset (brainstem pathways $\approx 3.0$\,ms,
reciprocal and motor pathways $\approx 1.8$\,ms, commissural
$\approx 2.2$\,ms) with $0.2$\,ms across-connection jitter. The
muscle$\to$Ia spindle/GTO feedback is implemented as a rate-coded
transduction (force and length $\to$ Ia firing rate), not a synaptic
projection, and so has no weight/delay entry.
Figure~\ref{fig:connectivity} shows the weight and delay distributions.


\begin{figure}[t]
  \centering
  \includegraphics[width=0.8\textwidth]{fig_connectivity.png}
  \caption{Distribution of synaptic weights and conduction delays along
  every projection (sensory-arm configuration: frozen BS$\to$RG, plastic
  Ia$\to$RG). \textbf{(a)} Weight histograms for the plastic / learned
  projections (STDP + lognormal init). \textbf{(b)} Synaptic weight per
  projection (mean $\pm$ per-connection s.d.; green excitatory, red
  inhibitory) --- all projections are heterogeneous (lognormal, CV$\approx
  0.5$), with the asymmetric reciprocal inhibition InF$\to$RG-E $=-48$\,pA
  versus InE$\to$RG-F $=-8$\,pA preserving the Zhang~6:1 ratio in the mean.
  \textbf{(c)} Conduction~+~synaptic delay per projection (mean $\pm$
  across-connection s.d.), following the rat \texttt{length\_velocity}
  preset with $0.2$\,ms jitter.}
  \label{fig:connectivity}
\end{figure}

\subsection{Brainstem drive}\label{sec:methods:bs}
Reticulospinal input is modelled as a tonic Poisson process at
$BS_{\text{regular}}=60$\,Hz with $\sigma=0.25$\,Hz noise, identical for
both legs. The 60\,Hz value sits in the
middle of the rat reticulospinal range
(20--80\,Hz)~\cite{drew1986,brocard2010}. In contrast to~\cite{zhang2022},
who vary tonic drive $\alpha$ to vary gait speed, we hold BS tonic and
pace the rhythm via sensory feedback (see \autoref{sec:methods:paced}).

\subsection{Spike-timing-dependent plasticity}\label{sec:methods:stdp}
Two projections --- cutaneous CUT$\to$RG-E and brainstem
BS$\to$RG-\{E,F\} --- carry multiplicative STDP synapses (NEST
\verb|stdp_synapse|; $\tau_+=20$\,ms, default $\lambda=10^{-3}$,
$\alpha=0.95$, $\mu_+=\mu_-=0.4$, $w_{\max}=120$\,pA for CUT and
$w_{\max}=30$\,pA for BS)~\cite{biPoo1998,morrison2007}. Initial
weights are drawn from a lognormal distribution with mean $\mu$ and
coefficient of variation CV. We swept $(\mu,\text{CV})$ across ten
diagnostic points (the ten $(\mu,\mathrm{CV})$ initialisation points) spanning four orders of
magnitude to verify that the post-training weight distribution
converges to the same plateau regardless of initialisation (see
Section~\ref{sec:results:weights}).

\paragraph{Extensor-biased plasticity and stance/swing asymmetry.}
The plastic and paced drives are deliberately extensor-weighted, mirroring
the stance/swing organisation of locomotor control. The cutaneous
afferent --- the only projection with the large ($w_{\max}=120$\,pA)
plasticity ceiling --- targets the extensor (CUT$\to$RG-E) and has no
flexor counterpart, reflecting that paw-contact cutaneous feedback fires
during weight-bearing stance and reinforces the load-bearing
extensor~\cite{pearson1995,rossignol2006}. Likewise the paced
heel$\to$toe Ia-E ramp clocks the extensor during stance, whereas the
flexor receives no external pacing and is generated intrinsically ---
by the persistent-sodium-like bursting of the RG-F units and by release
from reciprocal inhibition during swing. As a consequence the extensor
half-centre is sensory-driven and tightly clocked while the flexor is
centrally generated and intrinsically more variable; and because the
dominant plastic projection is extensor-only, increasing $\lambda$
strengthens the descending drive asymmetrically, in favour of the
extensor (analysed in Section~\ref{sec:results:force}).

\paragraph{Learning-rate variation.}
For the speed-by-STDP and graded-ablation matrices
(\autoref{sec:methods:protocols}), the STDP learning rate $\lambda$ was
varied across three values spanning two decades at and below the
cortical bio-plausible STDP rate~\cite{biPoo1998,morrison2007}:
$\lambda \in \{1{\times}10^{-5},\,1{\times}10^{-4},\,1{\times}10^{-3}\}$.
We refer to these as \emph{slow}, \emph{intermediate}, and
\emph{baseline} STDP, respectively. With a 120\,s run the slow rate does
not fully converge (time constant $\approx 250$\,s), which lets us probe
the regime of weak, still-developing descending weights. The parameter
is recorded to the simulation output for every run.

\subsection{Sensory feedback and Ia rate coding}\label{sec:methods:ia}
Ia afferent firing rate $r_{\text{Ia}}$ is updated every 100\,ms as a
linear combination of muscle force $F$ and stretch $\Delta\ell$:
\begin{equation}
  r_{\text{Ia}} = r_0 + k_F \cdot F + k_S \cdot \max(0, \ell - L_0),
  \quad \text{clamped to } [0, r_{\max}],
  \label{eq:ia}
\end{equation}
with $r_0=10$\,Hz, $k_F=6$\,Hz/F, $k_S=250$\,Hz/L,
$r_{\max}=500$\,Hz~\cite{loeb1981,prochazka1999}. This rate drives the
classical Ia-INT$\to$antagonist motor-pool
pathway~\cite{jankowska1992} and the closed-loop Ia$\to$recip-IN pathway
introduced in this work.

\paragraph{Graded sensory ablation (Courtine/Lavrov paradigm).}
To emulate the partial-weight-bearing and limb-unloaded preparations
used in rat SCI rehabilitation
experiments~\cite{lavrov2008,courtine2009,edgerton2008}, we expose a
multiplicative Ia-feedback gain $G_{\text{Ia}} \in [0, 1]$, applied
after Eq.~\ref{eq:ia}:
$\tilde{r}_{\text{Ia}} = G_{\text{Ia}}\cdot r_{\text{Ia}}$. The three
operating regimes used in Section~\ref{sec:results:epidural} are:
$G_{\text{Ia}}=1.0$ (\emph{baseline}, full weight-bearing);
$G_{\text{Ia}}=0.5$ (\emph{toe stepping}, partial weight); and
$G_{\text{Ia}}=0.1$ (\emph{air stepping}, limb unloaded, residual
spindle background only). The classical Ia-INT$\to$antagonist pathway
inherits the same gain, so the manipulation reduces \emph{all} Ia
output uniformly.

\subsection{Activation and force dynamics}\label{sec:methods:muscle}
The muscle relay rate $r_{\text{mus}}$ is mapped to activation $a$
through a saturating nonlinearity gated by RG-rate. The gate is a smooth
logistic (sigmoid) of the normalised RG population rate --- a
bio-plausible soft saturation in place of a hard clip:
\begin{align}
  a_{\text{raw}} &= a_{\max} \, \bigl(1 - \exp(-k_a \, r_{\text{mus}})\bigr), \label{eq:araw}\\
  d &= \sigma\!\Bigl(k_g\bigl(r_{\text{RG}}/r_{\text{ref}} - x_0\bigr)\Bigr),
       \qquad \sigma(z)=\frac{1}{1+e^{-z}}, \label{eq:gate}\\
  a_{\text{tgt}} &= a_{\text{raw}} \cdot d, \label{eq:atgt}\\
  a(t+\Delta t) &= a(t) + \bigl(1 - e^{-\Delta t/\tau_a}\bigr)\bigl(a_{\text{tgt}} - a(t)\bigr).
  \label{eq:alp}
\end{align}
with gate steepness $k_g=8$ and mid-point $x_0=0.6$. The logistic
saturates smoothly at $1$ and suppresses the inter-burst trough
($d\approx0.06$ at $r_{\text{RG}}/r_{\text{ref}}\approx0.25$ versus
$d\approx0.83$ at a burst peak), so the gated target $a_{\text{tgt}}$ is
bounded in $[0,a_{\max})$ by construction --- without the hard clip used
in earlier versions of the model.
Force follows an identical first-order low-pass driven by a saturating
function of activation:
\begin{align}
  F_{\text{tgt}} &= F_{\max} \, \bigl(1 - \exp(-k_F \, a)\bigr),\\
  F(t+\Delta t) &= F(t) + \bigl(1 - e^{-\Delta t/\tau_F}\bigr)\bigl(F_{\text{tgt}} - F(t)\bigr).
\end{align}
Under the paced-gait protocol (\autoref{sec:methods:paced}) the time
constants are overridden: $\tau_a=40$\,ms, $\tau_F=80$\,ms. The latter
sits within the rat fast-twitch fibre rise-time
range~\cite{zajac1989,winters1995}.

\paragraph{Force units and SI calibration.}
Force is reported in arbitrary units (a.u.). The dimensionless ceiling
$F_{\max}=25$\,a.u.\ in Eq.~(6) was chosen so the activation--force
cascade operates in a numerically convenient range; absolute Newton
calibration is not intrinsic to the present model and is conventional
in CPG-modelling reports~\cite{rybak2006,danner2017,zhang2022}. A
first-order linear rescaling maps a.u.\ to Newtons by anchoring on
\emph{in vivo} rat hindlimb data: the peak isometric tetanic force of
rat medial gastrocnemius is approximately
$5$--$7$\,N~\cite{roy1991,burkholderLieber2001}. Treating our reported
peak Force-E of $17.4$\,a.u.\ as representing peak ankle-extensor
tetanic force gives a conversion of $1$\,a.u.\ $\approx$
$0.30$--$0.40$\,N. Nothing in the model dynamics depends on this
choice of units; we use a.u.\ throughout for transparency.

\subsection{Paced gait protocol}\label{sec:methods:paced}
Each step period (default 1000\,ms) is divided
into two 500\,ms half-cycles. During the stance half of leg L, the L
cutaneous generator fires at $CUT_{\text{on}}=100$\,Hz, and three
sequential Ia-E sub-groups fire at 60, 80, and 100\,Hz over consecutive
$\sim$167\,ms windows, emulating heel$\to$mid$\to$toe pressure transfer
during stance~\cite{loebDuysens2018}. During swing, all CUT and Ia-E
drive is silenced. The contralateral leg R follows the 180\,$^{\circ}$
mirror schedule (trot pattern). Reciprocal inhibition and the RGF
intrinsic-bursting dynamics maintain the swing-leg flexor burst without
external pacing.

\subsection{Simulation setup}\label{sec:methods:sim}
All simulations were performed with NEST~3.9~\cite{nest2002} using a
kernel resolution of 0.2\,ms, simulation chunks of 100\,ms, and 64
OpenMP threads per task on a single AMD EPYC node (MN5 system,
Barcelona Supercomputing Center). Initial-weight sweeps were run as
10-task SLURM job arrays; ablation and speed-sweep experiments used
5-task arrays. A 30\,s simulation completes in $\sim$30\,min wall-clock;
a 120\,s simulation in $\sim$2\,h.

\subsection{Experimental protocols}\label{sec:methods:protocols}
Three classes of in-silico experiments are reported in this paper. Each
is encoded as a stand-alone shell script that submits a single
multi-task SLURM array and writes one HDF5 file per array task.
Pseudocode for all three is given in the Supplementary material (Algorithms~\ref{alg:main}--\ref{alg:ablgrad}).

\paragraph{Main experiment: long-term self-organisation.}
A 10-point sweep over the initial-weight distribution $(\mu,\mathrm{CV})$
for the STDP synapses, run as a SLURM job array. Each task simulates
\SI{120}{\second} of paced locomotion at $\mu, \mathrm{CV}$ drawn from
the ten $(\mu,\mathrm{CV})$ initialisation points. Used in Section~\ref{sec:results:weights}.


\paragraph{Phase A --- speed-by-STDP matrix.}
Nine 120\,s runs spanning three walking speeds (rat locomotor range,
\cite{lavrov2008,courtine2009}) crossed with three bio-plausible STDP
learning rates. Used in Section~\ref{sec:results:force}.


\paragraph{Phase B --- graded sensory ablation by STDP matrix
(Courtine/Lavrov SCI paradigm).}
Nine 120\,s runs at the medium-walk anchor ($P=520$\,ms) crossing three
Ia feedback gains with three STDP rates. Used in
Section~\ref{sec:results:epidural}.


\paragraph{Legacy ablation experiment (supplementary).}
For completeness we retain the earlier 5-condition ablation experiment
(\emph{baseline}, \emph{noIa}, \emph{symInh}, \emph{noComm},
\emph{noPaced}) that demonstrated which architectural components are
load-bearing. The Phase B graded ablation is the principal ablation
result of the paper; the binary ablation is reported in supplementary
Fig.~S1.

\noindent All three scripts (\texttt{run.sh}, \texttt{run\_ablation.sh},
\texttt{run\_speed.sh}) are publicly available in the repository. Each
single 30\,s task runs in $\sim$30\,min wall-clock on a 64-thread MN5
node; the 120\,s task takes $\sim$2\,h. Total compute for all three
experiments: $\sim$60 core-hours.

\subsection{Data analysis}\label{sec:methods:analysis}
Cycle period was computed as the Force-E peak-to-peak interval over the
last 20\,s, using a 5\,a.u.\ minimum peak height and a 300\,ms minimum
inter-peak gap. Counter-phase strength was reported as Pearson's
correlation between Force-E and Force-F over the same window. Peak
force amplitude was reported as the 95th percentile to avoid outlier
sensitivity. STDP convergence was reported as the per-projection mean
$\pm$ standard deviation of synaptic weights sampled every 1\,s. No
formal statistical tests are reported: simulations are deterministic with
a fixed seed (12345 + 10007$\times$idx for sweep tasks); reproducibility
is therefore exact, and uncertainty is captured by the $(\mu,\text{CV})$
initialisation sweep rather than by stochastic repetition.

% --- Parameter tables go in supplementary ---
